Layer-aware lookahead construction for 3D FPGA routing preserves source-relative delay and congestion costs, but its execution order can restrict parallelism both across samples and within a search. We expose these two granularities through context-preserving sample decoupling and whole-pair FP32 bucket waves. Private sample state and original-order reduction provide conditional native wire-map equivalence for STRICT CPU execution. ALT CPU/CUDA use a distinct wave selection rule, qualified separately through complete outputs and native routing. Matched scheduling controls identify cross-group parallelism, while wave profiles quantify concurrency versus repeated work. On new matched GEMM/BNN runs, CUDA achieves \GemmNewPipeProc\ and \BnnNewPipeProc\ complete-process speedups over eight-core Pipeline. A separate five-arm LSTM cohort gives $3.207\times$ construction and $1.094\times$ process speedups against Pipeline. Small-input slowdowns, representative-cost differences and cache applicability delimit these benefits.
